\documentclass[11pt]{article}

% ---------------------------------------------------------------------------
% Staff: release information.
\newcommand{\assigneddate}{Wednesday, October 21, 2026}
\newcommand{\duedate}{Wednesday, November 18, 2026}
\newcommand{\coderelease}{the starter-code folder distributed with this handout}
\newcommand{\staffsft}{\file{staff\_sft/} (in your starter-code folder)}
% Result thresholds = 0.9 x min staff reference accuracy over 2 runs (grading/thresholds.json).
\newcommand{\thrsft}{0.594}
\newcommand{\thrdpo}{0.617}
\newcommand{\thrgrpo}{0.677}
% ---------------------------------------------------------------------------

\usepackage[margin=1.1in]{geometry}
\usepackage{lmodern}
\usepackage[T1]{fontenc}
\usepackage{amsmath,amssymb}
\usepackage{graphicx}
\usepackage{booktabs}
\usepackage{array}
\usepackage{enumitem}
\usepackage{framed}
\usepackage{fancyhdr}
\usepackage[dvipsnames]{xcolor}
\usepackage{listings}
\usepackage[colorlinks=true,linkcolor=blue,urlcolor=blue,citecolor=blue]{hyperref}

\pagestyle{fancy}
\fancyhf{}
\renewcommand{\headrulewidth}{0pt}
\cfoot{Page \thepage}
\fancypagestyle{plain}{\fancyhf{}\renewcommand{\headrulewidth}{0pt}\cfoot{\thepage}}

\lstset{basicstyle=\ttfamily\small,columns=fullflexible,keepspaces=true,breaklines=true,
        xleftmargin=1.5em,frame=none,upquote=true}

% Inside solution boxes, figure environments become non-floating (a float inside a framed box is a LaTeX
% error), so \begin{figure}...\caption{...}\end{figure} works as usual.
\makeatletter
\newenvironment{solution}{\begin{framed}\noindent\textbf{Solution:}\par
  \renewenvironment{figure}[1][]{\par\medskip\centering\def\@captype{figure}}{\par\medskip}}{\end{framed}}
\makeatother
\setlength{\emergencystretch}{3em}
\newcommand{\pts}[1]{\textnormal{[#1~pts]}}
\newcommand{\file}[1]{\texttt{#1}}
\newcommand{\resultfigure}[3]{%
  \IfFileExists{#2}{%
    \begin{figure}
      \includegraphics[width=#1]{#2}
      \caption{#3}
    \end{figure}%
  }{\par\noindent\emph{[Pending figure]} \texttt{\detokenize{#2}}\par}%
}

\title{CSE595: Natural Language Processing\\Homework 3\\[0.3em]
\large Post-Training a Small Language Model to Reason}
% Students: put your name and uniqname here.
\author{[Kyungho (Matthew) Song] (\texttt{[songkh]})}
\date{Assigned date: \assigneddate{} --- Due date: \duedate}

\begin{document}
\maketitle

\noindent\textcolor{red}{\textbf{This assignment requires GPU computation on Great Lakes, a shared computing cluster at
UMich (about 2--3 A40 GPU-hours per student). You may need to wait in a queue to run your jobs, so please
start as early as possible!}}

\section*{Use of AI Tools}
While AI can be a valuable learning tool, it should not replace your own reasoning, critical thinking, and
problem-solving. You may use AI tools to help clarify concepts, understand algorithms, explore programming
techniques, and debug your code. However, you may not simply provide an assignment to an AI tool and ask it to
complete the work for you. You are expected to develop your own solutions and fully understand the work you
submit. You must \textbf{fully disclose how you used AI tools} in the ``Disclosure'' section of the assignment, including
the tools used and the nature of the assistance received.

\section*{Submission Guidelines}
\begin{enumerate}
  \item Please insert your student information (name and uniqname) in the \verb|\author| line of this \LaTeX{} file.
  \item Please insert your answers between each pair of \verb|\begin{solution}| and \verb|\end{solution}|.
  \item Get the starter code from \coderelease. Implement the \texttt{TODO}s in
        \file{posttrain/\{common, sft, preference, grpo, test\_time, gsm8k\}.py}. Do not edit the other provided
        files (see ``Grading'').
  \item We use \href{https://www.gradescope.com/courses/1384990}{Gradescope} for homework submission.
  \item We have two submission portals on Gradescope: \mbox{\texttt{Homework 3 - Coding}} and
        \mbox{\texttt{Homework 3 - Written}}.
  \item On Great Lakes, package your code and results with
        \begin{lstlisting}
python scripts/make_submission.py --uniqname <uniqname>\end{lstlisting}
        and submit \file{submission\_<uniqname>.zip} to \mbox{\texttt{Homework 3 - Coding}} (83 pts). The coding
        problems will be graded automatically. Checklist for the zip (the script reports anything missing):
        \begin{itemize}[nosep]
          \item \file{posttrain/\{common, sft, preference, grpo, test\_time, gsm8k\}.py};
          \item \file{results/\{sft, dpo, grpo\}/\{predictions.jsonl, metrics.json\}};
          \item \file{results/test\_time/test\_time.json};
          \item \file{results/gsm8k/\{A\_standard, B\_standard, A\_improved\}/predictions.jsonl}.
        \end{itemize}
  \item Please compile this file to \file{hw3.pdf} (e.g.\ in Overleaf) and submit the PDF to
        \mbox{\texttt{Homework 3 - Written}} (Problems 2--7, 17 pts, plus the required Disclosure section).
        {\color{red}Please make sure that each question is matched to the correct page on Gradescope. Each
        missing or mismatched question will result in a 1-point deduction.}
\end{enumerate}

\section*{Overview}
In this assignment you will take a small \emph{base} language model, Qwen2.5-0.5B~\cite{qwen25}, which cannot
follow instructions, and walk it through the modern post-training pipeline on a single GPU:
\begin{enumerate}
  \item \textbf{Supervised fine-tuning (SFT)} on reasoning traces, so the model learns the output format and a
        basic search strategy.
  \item \textbf{Preference optimization (DPO~\cite{dpo} or IPO~\cite{ipo})} on pairs of good and bad answers,
        labelled automatically by a verifier. This is RLHF~\cite{instructgpt} with the verifier standing in for
        the human labeller.
  \item \textbf{Reinforcement learning with a verifiable reward (GRPO~\cite{deepseekmath})}, where the model
        improves from its own rollouts.
  \item \textbf{Test-time scaling}, where you trade extra inference compute for accuracy with majority
        voting~\cite{selfconsistency}, confidence-based selection and pass@$k$~\cite{codex}.
  \item \textbf{Benchmarking}, where you evaluate two released instruct models on a standard math benchmark,
        GSM8K~\cite{gsm8k}, and measure how much the reported accuracy depends on the prompt.
\end{enumerate}
You implement the core math of every stage: masked log-likelihoods, the SFT loss, the DPO/IPO losses,
group-relative advantages, the clipped policy-gradient objective, and the test-time selection rules. We provide
the training loops, data, evaluation and Slurm scripts, and they call your functions.

\begin{center}
\begin{tabular}{@{}llr@{}}
\toprule
component & how it is graded & points \\
\midrule
Unit tests (Problems 1--5) & autograded against hidden test cases & 50 \\
SFT / DPO / GRPO results & autograded from your \file{predictions.jsonl} & 10 / 6 / 10 \\
Test-time artifact & autograded consistency check of \file{test\_time.json} & 4 \\
GSM8K (Problem 7) & autograded from your GSM8K predictions & 3 \\
Write-up (Problems 2--7) & graded by hand from your solution boxes & 17 \\
\midrule
Total & & 100 \\
\bottomrule
\end{tabular}
\end{center}

\section*{The Task: Countdown}
Given a few integers and a target, write an arithmetic expression that uses \textbf{each number exactly once}
with \texttt{+ - * /} and parentheses and equals the target~\cite{streamofsearch}. For example:
\begin{lstlisting}
Numbers [44, 19, 35], target 98   ->   <answer>(44 + 19) + 35</answer>\end{lstlisting}
A deterministic verifier (\file{posttrain/countdown.py}) checks answers exactly using rational arithmetic, so
there is no LLM judge and no human labelling. The reward is \textbf{1.0} if the answer is correct, \textbf{0.1}
if an \texttt{<answer>} block exists but is wrong, and \textbf{0.0} otherwise. All problems use 3 or 4 numbers.
The data is in \file{data/}:
\begin{center}
\begin{tabular}{@{}lrp{6.4cm}@{}}
\toprule
file & size & use \\
\midrule
\file{sft\_train.jsonl} / \file{sft\_val.jsonl} & 8000 / 200 & prompt + synthetic search trace (a few failed tries, then a correct one) \\
\file{countdown\_train.jsonl} & 20000 & prompts for preference-pair sampling and RL \\
\file{countdown\_val.jsonl} & 200 & monitoring during RL \\
\file{pref\_pairs.jsonl} & 1851 & preference pairs built from the staff SFT model (Problem 3) \\
\file{countdown\_test.jsonl} & 500 & \textbf{graded; never train or tune on it} \\
\bottomrule
\end{tabular}
\end{center}

\section*{Getting Started}
\textbf{Do all the implementation work on your laptop or a Great Lakes login node first.} Every unit test runs
on a CPU in seconds:
\begin{lstlisting}
pip install torch pytest
python -m pytest tests/                   # all fail at first: NotImplementedError
python -m pytest tests/test_sft.py -x -q  # one file, stop at the first failure\end{lstlisting}
On Great Lakes, set up the environment once from a login node, then submit jobs:
\begin{lstlisting}
bash slurm/setup_env.sh    # creates a venv in your scratch space (~10 min)
sbatch slurm/sft.sbatch    # then dpo, grpo, grpo_ablation, test_time, gsm8k
squeue -u $USER            # monitor; logs go to logs/\end{lstlisting}
Before a full run, smoke-test on a GPU with small sizes, for example \texttt{-{}-max-examples 256 -{}-epochs 1}
(SFT), \texttt{-{}-steps 3 -{}-prompts-per-step 8} (GRPO) and \texttt{-{}-limit 50 -{}-backend hf} (evaluation). \textbf{Graded evaluation must use
the full test set.} If a test fails, read the assertion message: it names the quantity and shows the expected
and actual values.

\paragraph{Plots for the write-up.} You do not need to write plotting code. After each stage, run
\begin{lstlisting}
python scripts/plot_results.py   # CPU only; run it on a login node\end{lstlisting}
It reads \file{results/} and saves every figure the write-up asks for (the file names are given in each
problem below; stages that have not run yet are skipped), and it prints a table of accuracy and format rate for
every finished run. Recompile this document with the generated figures at the listed paths to include them.

% =====================================================================================
\section*{Notation and Conventions}
Every training stage in this assignment works on the same kind of data: a \emph{prompt} $x$ (the Countdown
question) followed by a \emph{response} $y$ (the model's answer). The symbols below are used throughout the
handout and the docstrings.
\begin{center}
\begin{tabular}{@{}l>{\raggedright\arraybackslash}p{11.2cm}@{}}
\toprule
symbol & meaning \\
\midrule
$B$ & batch size: the number of (prompt, response) sequences processed together. \\
$L$ & length, in tokens, of the longest prompt+response sequence in the batch; shorter sequences are
      \emph{padded} on the right with a pad token up to length $L$. \\
$T$ & $= L-1$, the number of next-token predictions per sequence (see Problem 1, item 1). \\
$V$ & vocabulary size (151{,}936 for Qwen2.5). \texttt{model(input\_ids).logits} has shape $(B, T, V)$. \\
$\pi_\theta$ & the model being trained (the \emph{policy}), with parameters $\theta$. \\
$\ell_{b,t}$ & $\log\pi_\theta(\text{token } t{+}1 \text{ of sequence } b \mid \text{tokens } 1..t)$: the
      per-token log-probability of the next token, a $(B, T)$ tensor called \texttt{log\_probs} or
      \texttt{policy\_log\_probs} in the code. \\
$m_{b,t}$ & the \emph{response mask}: 1 if the token being predicted at position $t$ belongs to the response,
      0 if it belongs to the prompt or is padding. Losses only count positions with $m_{b,t}=1$, because we want
      to train the model to \emph{write answers}, not to reproduce the question. \\
$K$ & number of gradient-accumulation steps (\texttt{gradient\_accumulation\_steps}). A batch that does not fit
      in GPU memory is split into $K$ \emph{micro-batches}; each micro-batch calls \texttt{backward()} on its loss
      divided by $K$, so the gradients that add up in \texttt{.grad} equal the gradient of the batch-average
      loss. The optimizer then takes one step. \\
\bottomrule
\end{tabular}
\end{center}
\emph{Greedy decoding} always picks the most likely next token (temperature 0), so it is deterministic;
\emph{sampling} at temperature 1 draws each token at random from the model's distribution. \emph{Accuracy} is the fraction of test
problems the verifier marks correct; \emph{format rate} is the fraction of responses that contain an
\texttt{<answer>...</answer>} block at all, correct or not.

% =====================================================================================
\section*{Problem 1: Shared Utilities \pts{12}}
All functions are in \file{posttrain/common.py} and tested by \file{tests/test\_common.py} (12 pts, autograded).
These are the building blocks that SFT, DPO and GRPO all reuse: turn text into tensors, ask the model how likely
each response token is, and average over the response tokens only. Read each docstring carefully: the tests check
exact shapes and conventions.
\begin{enumerate}
  \item (Programming) Please implement \texttt{tokenize\_prompt\_and\_output()}. Tokenize each prompt and each
        output \emph{separately}, concatenate them, and right-pad the batch with the pad token to get a $(B, L)$
        tensor $\texttt{full}$. Return $\texttt{input\_ids} = \texttt{full}[:, :\!-1]$,
        $\texttt{labels} = \texttt{full}[:, 1:]$ and a boolean \texttt{response\_mask} over \texttt{labels} that
        is true exactly where the label token belongs to the output (not the prompt, not padding).

        \emph{Why the shift?} A language model reads tokens $1..t$ and predicts token $t{+}1$, so the target
        (\texttt{labels}) at position $t$ is the input shifted left by one. This is called \emph{teacher forcing}:
        one forward pass scores every response token at once. Example with $B=2$, where sequence~1 has prompt
        tokens \texttt{p1 p2 p3} and output tokens \texttt{o1 o2}, and sequence~2 has prompt \texttt{q1 q2} and
        output \texttt{r1} (so $L=5$ and $T=4$):
\begin{center}
\begin{tabular}{@{}l|cccc|cccc@{}}
\toprule
 & \multicolumn{4}{c|}{sequence 1} & \multicolumn{4}{c}{sequence 2} \\
\midrule
\texttt{full}           & \multicolumn{4}{c|}{\texttt{p1 p2 p3 o1 o2}} & \multicolumn{4}{c}{\texttt{q1 q2 r1 PAD PAD}} \\
\texttt{input\_ids}     & \texttt{p1} & \texttt{p2} & \texttt{p3} & \texttt{o1} & \texttt{q1} & \texttt{q2} & \texttt{r1} & \texttt{PAD} \\
\texttt{labels}         & \texttt{p2} & \texttt{p3} & \texttt{o1} & \texttt{o2} & \texttt{q2} & \texttt{r1} & \texttt{PAD} & \texttt{PAD} \\
\texttt{response\_mask} & F & F & T & T & F & T & F & F \\
\bottomrule
\end{tabular}
\end{center}
  \item (Programming) Please implement \texttt{compute\_entropy()}, the entropy
        $H(p) = -\sum_v p_v \log p_v$ (in nats, i.e.\ natural log) of $p=\mathrm{softmax}(\text{logits})$ over the
        last dimension. Entropy measures how spread out the next-token distribution is; GRPO logs it to show
        whether the model is becoming more deterministic.

        \emph{Numerical stability.} Logits can be large, e.g.\ around $10^3$ in absolute value (the hidden tests
        use such values). Computing $p_v = e^{z_v}/\sum_u e^{z_u}$ directly overflows: $e^{1000}$ is
        \texttt{inf} in floating point, and then $\log 0 = -\infty$ and $0\cdot(-\infty)$ give \texttt{nan}.
        Instead compute $\log p$ with \texttt{torch.log\_softmax} (which subtracts the maximum logit internally
        and never forms $e^{1000}$), then $H = -\sum_v e^{\log p_v}\log p_v$. Your function must return finite
        values for such inputs.
  \item (Programming) Please implement \texttt{get\_response\_log\_probs()}: for every position, the
        log-probability of the label token given the prefix, $\ell_{b,t} = \log \pi_\theta(\texttt{labels}[b,t]
        \mid \texttt{input\_ids}[b, 0..t])$, computed in float32 from \texttt{model(input\_ids).logits},
        optionally together with the per-token entropy. In words: take the log-softmax over the vocabulary at
        every position, then pick out the entry of the token that actually comes next (\texttt{torch.gather}).
        The result has shape $(B, T)$; masking happens later, in the loss.
  \item (Programming) Please implement \texttt{masked\_mean()} and \texttt{masked\_normalize()}: the mean, and
        the sum divided by a constant, of a tensor over the positions where a mask is nonzero, either over all
        elements or along one dimension. For example, with
        $\texttt{tensor} = [[1, 2, 3], [4, 5, 6]]$ and $\texttt{mask} = [[1, 1, 0], [0, 1, 1]]$:
        \texttt{masked\_mean} with \texttt{dim=None} is $(1+2+5+6)/4 = 3.5$, with \texttt{dim=1} it is
        $[1.5, 5.5]$, and \texttt{masked\_normalize} with constant $2$ and \texttt{dim=1} is $[1.5, 5.5]$
        (sums $3$ and $11$, divided by $2$). The losses in Problems 2--4 use these to average over response
        tokens only.
\end{enumerate}
\textbf{Hint.} The \texttt{response\_mask} convention (which positions count) is the most common source of bugs.
Draw a picture of one short example, like the table above, before you write code.

% =====================================================================================
\section*{Problem 2: Supervised Fine-Tuning \pts{18}}
\begin{enumerate}
  \item (Programming, 6 pts) Please implement the SFT loss in \file{posttrain/sft.py}, function
        \texttt{sft\_microbatch\_train\_step()}. SFT is ordinary next-token cross-entropy on the response
        tokens: we maximize the log-probability of the reference solution (here, a synthetic search trace that
        tries a few expressions and ends with a correct one). With the per-token log-probabilities $\ell_{b,t}$ and
        response mask $m_{b,t}$ of one micro-batch (see Notation), a normalizing constant $c$
        (\texttt{normalize\_constant}) and $K$ gradient-accumulation steps, the loss is
        \[
          \mathcal{L}_\text{SFT} = -\frac{1}{K}\cdot\frac{1}{B}\sum_{b=1}^{B}\frac{1}{c}\sum_{t} m_{b,t}\,\ell_{b,t}.
        \]
        The minus sign turns ``maximize log-likelihood'' into a loss to minimize. The training script passes
        $c$ = the average number of response tokens per sequence in the micro-batch, which makes the loss a
        per-token average. Call \texttt{backward()} on it inside the function.
  \item (Experiment, 10 pts) Train and evaluate. The job first evaluates the \emph{base} model (Qwen2.5-0.5B before
        any training) once for comparison, then runs SFT for 2 epochs (2 passes over the 8000 training examples)
        and evaluates your model with greedy decoding on the 500 test problems:
\begin{lstlisting}
sbatch slurm/sft.sbatch   # -> results/base/ and results/sft/\end{lstlisting}
        Full credit requires a greedy test accuracy of at least \thrsft; below that, credit is proportional to
        accuracy.
  \item (Write-up, 2 pts) Include the training and validation loss plot, \file{results/sft/loss.png}.
        Report greedy accuracy and format rate for the base model and for your SFT model (\file{plot\_results.py}
        prints both; they come from the \file{metrics.json} files). Why do we mask prompt tokens out of the loss, and what would change
        if we did not?
\end{enumerate}
\begin{solution}
\resultfigure{0.6\linewidth}{results/sft/loss.png}{SFT training and validation loss.}

The base model reached [base accuracy] greedy accuracy with a format
rate of [base format rate]. After SFT, greedy accuracy was [SFT accuracy]
and the format rate was [SFT format rate].

Masking the prompt makes the objective the conditional log-likelihood of
the response given the question. The model still reads the prompt when
predicting response tokens. Including prompt tokens would also reward
predicting the question itself, and longer prompts could take up a larger
share of the loss.
\end{solution}

% =====================================================================================
\section*{Problem 3: Preference Optimization \pts{19}}
Preference optimization learns from \emph{pairs}: for the same prompt $x$, a \emph{chosen} (better) response
$y_w$ and a \emph{rejected} (worse) response $y_l$. Instead of imitating one reference answer as in SFT, the model
is pushed to make $y_w$ more likely \emph{relative to} $y_l$. Write $\pi_\theta$ for the policy being trained,
$\pi_\text{ref}$ for a frozen copy of the SFT model (the \emph{reference policy}, which is never updated), and
$\log\pi(y\mid x)$ for the log-probability of a whole response, i.e.\ the sum of its per-token
log-probabilities. Let
\[
  h = \bigl[\log\pi_\theta(y_w\mid x) - \log\pi_\text{ref}(y_w\mid x)\bigr]
    - \bigl[\log\pi_\theta(y_l\mid x) - \log\pi_\text{ref}(y_l\mid x)\bigr].
\]
Each bracket measures how much more likely the policy makes a response than the reference does; DPO calls
$\beta$ times this bracket the \emph{implicit reward} of the response. So $h>0$ means the policy has moved
towards the chosen response and away from the rejected one. The hyper-parameter $\beta>0$ (default 0.1)
controls how far the policy may drift from the reference: smaller $\beta$ allows larger changes.
\paragraph{Staff SFT checkpoint.} If your SFT model is weak (for example, below the SFT threshold), you may start
Problems 3--5 from the staff SFT checkpoint,
\begin{center}\staffsft,\end{center}
by submitting jobs with \texttt{USE\_STAFF\_SFT=1}, e.g.\ \texttt{USE\_STAFF\_SFT=1 sbatch slurm/dpo.sbatch}. Your SFT result is still graded on your own SFT model, and
DPO/GRPO results are graded the same way either way; just say in your write-up which starting point you used.
\begin{enumerate}
  \item (Programming, 10 pts for items 1--2) Please implement \texttt{sequence\_log\_probs()} in
        \file{posttrain/preference.py}: the sum of token log-probabilities over the response tokens of each sequence.
  \item (Programming) Please implement \texttt{preference\_loss()}, with
        $\mathcal{L}_\text{DPO} = -\log\sigma(\beta h)$, where $\sigma(z) = 1/(1+e^{-z})$ is the sigmoid, and
        $\mathcal{L}_\text{IPO} = \bigl(h - \tfrac{1}{2\beta}\bigr)^2$, both averaged over the batch. DPO keeps
        pushing $h$ up; IPO pulls $h$ towards the fixed target $1/(2\beta)$, which limits over-fitting. Also return
        the four metrics described in the docstring: chosen reward $\beta[\log\pi_\theta(y_w)-\log\pi_\text{ref}(y_w)]$,
        rejected reward (the same for $y_l$), reward margin $\beta h$, and reward accuracy (the fraction of
        pairs with $h>0$). Use \texttt{torch.nn.functional.logsigmoid} rather than \texttt{log(sigmoid(...))}, for
        the same numerical reason as in Problem 1.
  \item (Experiment, 6 pts) We provide \file{data/pref\_pairs.jsonl}. To build it, the staff decoded 7000 training
        prompts greedily with the staff SFT model; it got 2378 of them wrong. For each of those, the staff sampled 8
        responses ($T=1$) and labelled them with the verifier. Whenever at least one sample was correct, a correct
        sample became the chosen response and the model's own wrong \emph{greedy} completion became the rejected
        response. Run DPO against a frozen copy of your SFT model as the reference policy:
\begin{lstlisting}
sbatch slurm/dpo.sbatch   # DPO + evaluation -> results/dpo/\end{lstlisting}
        Full credit requires a greedy test accuracy of at least \thrdpo. Keep the default learning rate
        ($3\times10^{-7}$): with these pairs, staff runs at $7\times10^{-7}$ and above collapsed.
        (Optional: delete or rename \file{data/pref\_pairs.jsonl} and the job will build the pairs from your own
        SFT model, which adds about 15 minutes.)
  \item (Write-up, 3 pts) Include the plot of validation reward accuracy and reward margin over training
        (\file{results/dpo/rewards.png}), and compare your DPO test accuracy to your SFT test accuracy. In a staff experiment, pairs whose rejected response was a wrong \emph{sampled} response
        instead of the wrong greedy completion did not improve greedy accuracy at all. Look at the rejected responses
        in \file{data/pref\_pairs.jsonl} and at your SFT model's wrong greedy completions in
        \file{results/sft/predictions.jsonl}, and explain why the choice of rejected response matters. In 3--5
        sentences, explain why this pipeline counts as ``RLHF'' even though no human labelled anything, and what role
        $\beta$ plays. Finally, only the 2378 greedy failures were used, and 527 of them were then discarded because
        none of the 8 samples was correct. What does that imply about which problems DPO can learn from?
\end{enumerate}
\begin{solution}
\resultfigure{0.95\linewidth}{results/dpo/rewards.png}{DPO validation reward accuracy and margin.}

DPO started from [SFT checkpoint used]. Its greedy test accuracy was
[DPO accuracy], compared with [SFT accuracy] for my SFT model.

In the supplied preference data, 1,790 of the 1,851 rejected completions
lack a complete answer block. For example, in \texttt{train-00001},
the rejected completion repeatedly evaluates variations of the same
expression to 39 for a target of 83. My SFT completion for [test ID]
shows [observed behavior]. A rejected greedy completion targets an
error the decoder actually returns at evaluation time. A wrong sampled
response can have low probability already, so reducing its probability
may leave the greedy result unchanged.

The verifier supplies the chosen/rejected judgments that human raters
would supply in an RLHF pipeline. DPO uses these judgments to increase
the relative likelihood of the chosen response against a frozen SFT
reference. The parameter $\beta$ scales the policy-reference log-ratio
margin and controls the strength of the reference constraint.

The 1,851 pairs cover prompts where greedy decoding failed and at least
one of eight samples succeeded. Prompts already solved greedily and
prompts with no correct sample among the eight provide no direct
training signal through these pairs.
\end{solution}

% =====================================================================================
\section*{Problem 4: RL with Verifiable Rewards \pts{28}}
In RL the model learns from its own attempts. One GRPO \emph{step}: take 32 training prompts, sample $G=8$
responses for each (a \emph{group}) at temperature 1, score every response with the verifier (reward 1.0 /
0.1 / 0.0, see ``The Task''), and update the model so that responses that scored better than the rest of their
group become more likely and worse ones less likely. No reference answer and no value network are needed: the
other samples for the same prompt serve as the baseline.

\paragraph{Advantages.} For a prompt with rewards $r_1,\dots,r_G$, the \emph{advantage} of response $i$ says how
much better it was than expected:
$A_i = (r_i - b_i)/(\mathrm{std}(r) + \delta)$ with either the group mean $b_i = \frac{1}{G}\sum_j r_j$ (GRPO)
or the leave-one-out mean $b_i = \frac{1}{G-1}\sum_{j\neq i} r_j$ (RLOO~\cite{rloo}). Dividing by the
\emph{population} standard deviation of the group's rewards (divide by $G$, not $G-1$) is optional;
$\delta = 10^{-6}$ (\texttt{eps} in the code) only avoids division by zero. Example: rewards $(1, 0, 0, 0)$ give
mean baseline $0.25$, so before normalization the advantages are $(0.75, -0.25, -0.25, -0.25)$. If all rewards
in a group are equal, every advantage is 0 and that prompt contributes no gradient.

\paragraph{Policy loss.} Each token of response $i$ gets the advantage $A_i$ of its whole response. Vanilla
policy gradient minimizes $-A_i\,\log\pi_\theta(y_t\mid\cdot)$. GRPO uses the PPO-style clipped
objective~\cite{ppo}: let $\pi_{\theta_\text{old}}$ be the model \emph{as it was when the responses were
sampled}, and $\rho_t = \pi_\theta(y_t\mid\cdot)/\pi_{\theta_\text{old}}(y_t\mid\cdot)$ the probability ratio
(computed as $\exp(\log\pi_\theta - \log\pi_{\theta_\text{old}})$). The per-token loss is
$-\min\bigl(\rho_t A,\ \mathrm{clip}(\rho_t, 1-\varepsilon, 1+\varepsilon)A\bigr)$ with
$\varepsilon = 0.2$ (\texttt{clip\_eps}), which removes the incentive to move a token's probability more than
$\pm 20\%$ away from $\pi_{\theta_\text{old}}$ in one update. A token is \emph{clipped} when the clipped term
is strictly smaller than the unclipped one; \texttt{clip\_fraction} is the fraction of response tokens that
are clipped.
\begin{enumerate}
  \item (Programming, 14 pts for items 1--4) Please implement \texttt{compute\_group\_advantages()} in
        \file{posttrain/grpo.py}, with the
        \texttt{mean} and \texttt{loo} baselines and optional standard-deviation normalization.
  \item (Programming) Please implement \texttt{per\_token\_policy\_loss()}, as vanilla policy gradient
        (\texttt{pg}, loss $-A\log\pi_\theta$) or the clipped objective above (\texttt{grpo\_clip}), also returning
        the mask of tokens where the clipped term is active.
  \item (Programming) Please implement \texttt{aggregate\_loss()}, which turns the $(B,T)$ per-token loss into
        one number, with three reductions over response tokens (they differ in how much weight long responses get):
        \texttt{seq\_mean} (average within each sequence, then over the batch; original GRPO),
        \texttt{token\_mean} (average over all response tokens; DAPO~\cite{dapo}) and \texttt{constant} (sum
        divided by batch size times a constant; Dr.~GRPO~\cite{drgrpo}).
  \item (Programming) Please implement \texttt{grpo\_microbatch\_train\_step()}: per-token loss, aggregation,
        division by the number of gradient-accumulation steps $K$, and \texttt{backward()}.
  \item (Experiment, 10 pts) Train with GRPO from your SFT model (100 steps $\times$ 32 prompts $\times$ 8 samples) and
        evaluate:
\begin{lstlisting}
sbatch slurm/grpo.sbatch   # -> results/grpo/\end{lstlisting}
        \file{results/grpo/train\_log.jsonl} has one line per step with the mean reward of the sampled responses,
        training accuracy, validation accuracy (greedy, on \file{countdown\_val.jsonl}, every 20 steps), mean
        token entropy (Problem 1), mean response length in characters (\texttt{response\_chars}),
        \texttt{clip\_fraction}, and
        \texttt{frac\_groups\_all\_same}, the fraction of groups in which all 8 samples got the same reward. Full
        credit requires a greedy test accuracy of at least \thrgrpo.
  \item (Experiment, part of the write-up points) Run \textbf{one ablation} for 40 steps, choosing one of
        \begin{itemize}[nosep]
          \item \texttt{-{}-baseline loo -{}-no-std-normalization} (RLOO),
          \item \texttt{-{}-aggregation constant -{}-no-std-normalization} (Dr.~GRPO),
          \item \texttt{-{}-epochs-per-rollout-batch 2} (two optimizer steps per batch of rollouts).
        \end{itemize}
        Pass the flags through \texttt{ABLATION}; the results go to \file{results/grpo\_ablation/}:
\begin{lstlisting}
ABLATION="--baseline loo --no-std-normalization" \
    sbatch slurm/grpo_ablation.sbatch\end{lstlisting}
  \item (Write-up, 4 pts) Include the plot of reward, validation accuracy, entropy and response length against
        step (\file{results/grpo/curves.png}; it also shows \texttt{frac\_groups\_all\_same} and
        \texttt{clip\_fraction}). Compare your ablation to the first 40 steps of the default run
        (\file{results/grpo\_ablation/compare.png}) and interpret the difference. What does
        \texttt{frac\_groups\_all\_same} tell you about the learning signal, and how would you expect it to change
        for problems that are too easy or too hard? Why is \texttt{clip\_fraction} exactly 0 throughout the
        default run?
\end{enumerate}
\begin{solution}
\resultfigure{0.95\linewidth}{results/grpo/curves.png}{GRPO training curves (default run).}
\resultfigure{0.95\linewidth}{results/grpo_ablation/compare.png}{Ablation vs.\ the first 40 steps of the default run.}

The planned ablation uses a leave-one-out baseline without
standard-deviation normalization. At step 40, validation accuracy is
[default validation accuracy] for the default run and [ablation
validation accuracy] for the ablation. The comparison plot also shows
[observed difference in reward, entropy, or response length].

With a fixed group size $G$, the unnormalized leave-one-out advantage is
$G/(G-1)$ times the unnormalized mean-baseline advantage. Removing
standard-deviation normalization also changes the relative weight of
groups with different reward variance. [Interpret the observed curves
in light of these two changes.]

\texttt{frac\_groups\_all\_same} measures the fraction of prompts for which
all sampled responses receive the same reward. Their advantages are
zero, so they provide no relative policy-gradient signal. This fraction
can rise for very easy prompts when every sample succeeds and for very
hard prompts when every sample earns the same low reward.

In the default run, all microbatch losses are computed before the one
optimizer step for that rollout batch. The current and cached old
log-probabilities therefore agree when the objective is evaluated.
The probability ratio is 1, so the clipping condition never activates.
\end{solution}

% =====================================================================================
\section*{Problem 5: Test-Time Scaling \pts{16}}
Instead of a better model, we can spend more compute at inference: sample $k$ responses per problem and choose
one. With a perfect checker you would just pick a correct sample if there is one; the probability that at least
one of $k$ samples is correct is \textbf{pass@$k$}. At deployment there is no checker, so we also need
\emph{verifier-free} selection rules: \textbf{maj@$k$} (majority vote: return the most common answer among the
$k$ samples) and \textbf{conf@$k$} (return the sample the model is most confident in, measured by its mean
token log-probability). maj@$k$ and conf@$k$ are the accuracies of these rules.
\begin{enumerate}
  \item (Programming, 8 pts for items 1--2) In \file{posttrain/test\_time.py}, please implement
        \texttt{pass\_at\_k()}, the unbiased estimator~\cite{codex} $\text{pass@}k = 1 - \binom{n-c}{k}\big/\binom{n}{k}$ for $n$ samples of which $c$ are
        correct. We draw $n \ge k$ samples once and estimate pass@$k$ from them: $\binom{n-c}{k}/\binom{n}{k}$ is the
        probability that $k$ samples picked without replacement from the $n$ are all wrong. Example: $n=4$, $c=1$,
        $k=2$ gives $1 - \binom{3}{2}/\binom{4}{2} = 1 - 3/6 = 0.5$. It must be numerically stable for $n$ up to
        1000: $\binom{1000}{500}$ has about 300 digits and overflows a float, so rewrite the ratio as a product of
        $k$ terms $\prod_{i=0}^{k-1}\frac{n-c-i}{n-i}$ (or work with logarithms).
  \item (Programming) Please implement \texttt{majority\_vote()} (most common non-\texttt{None} answer, ties go to
        the earliest first occurrence) and \texttt{select\_by\_confidence()} (index of the highest score;
        \texttt{None}/NaN count as $-\infty$; ties go to the smallest index).
  \item (Experiment, 4 pts) Sample 16 responses for each of 200 test problems from your GRPO model and from your SFT
        model, and compute \textbf{pass@$k$} (oracle upper bound), \textbf{maj@$k$} (vote over the \emph{value} of
        each answer) and \textbf{conf@$k$} (pick the sample with the highest mean token log-probability) for
        $k\in\{1,2,4,8,16\}$:
\begin{lstlisting}
sbatch slurm/test_time.sbatch   # -> results/test_time{,_sft}/\end{lstlisting}
        GRPO results go to \file{results/test\_time/} and SFT results to \file{results/test\_time\_sft/}.
        The 4 artifact points require a well-formed \file{results/test\_time/test\_time.json} with consistent
        curves (pass@$k$ non-decreasing in $k$, and maj@$k$ and conf@$k$ not above pass@$k$).
  \item (Write-up, 4 pts) Include the two plots (\file{results/test\_time/test\_time.png} for GRPO and
        \file{results/test\_time\_sft/test\_time.png} for SFT; \file{test\_time\_scaling.py} saves them). Explain the gap between pass@$k$ and the two verifier-free
        methods. Why is plain majority voting over answer \emph{strings} a poor fit for Countdown, and does voting
        over \emph{values} fix it? Using your numbers, is it better to sample 16 times from the SFT model or once
        from the GRPO model?
\end{enumerate}
\begin{solution}
\resultfigure{0.48\linewidth}{results/test_time/test_time.png}{Test-time scaling, GRPO model.}
\resultfigure{0.48\linewidth}{results/test_time_sft/test_time.png}{Test-time scaling, SFT model.}

Pass@$k$ measures how often a subset contains at least one correct
response. The oracle can always pick that response, while majority
voting and confidence selection can choose an incorrect response from
the same subset. String voting splits votes across different expressions
that evaluate to the same value. Voting by value pools their votes,
but an incorrect value can still win. An expression can also reach the
target while violating the requirement to use each given number once.

For SFT at $k=16$, pass@$k$, maj@$k$, and conf@$k$ were [three values].
For GRPO at $k=1$, the corresponding scores on the same problems were
[three values]. Under verifier-free selection, [state which setup was
more accurate and by how much]. With an oracle, [state how the
comparison changes, if it does].
\end{solution}

% =====================================================================================
\section*{Problem 6: Failure Analysis \pts{2}}
\begin{enumerate}
  \item (Write-up, 2 pts) Show one incorrect greedy completion from your best model (from
        \file{results/*/predictions.jsonl}) and diagnose it. Is it an arithmetic slip, a number-reuse violation, a
        formatting error, a search failure, or a repetition loop that runs out of the 384-token budget? How common
        is that failure type among your model's errors?
\end{enumerate}
\begin{solution}
The model with the highest greedy test accuracy was [model].
On [test ID], the task was to use [numbers] to reach [target].
The model produced the following completion.

\begin{lstlisting}
[Paste the actual completion.]
\end{lstlisting}

I classified this as [failure type] because [specific evidence from
the trace or final expression]. Using [classification rule], I found
[count] instances of this failure among [total] incorrect completions,
or [percentage]\% of the errors.
\end{solution}

% =====================================================================================
\section*{Problem 7: Closed-Ended Benchmarking \pts{5}}
So far you have measured accuracy on Countdown with our own verifier. Most published LM benchmarks work the same
way: a fixed set of \emph{closed-ended} questions, each with one correct answer, a fixed prompt, a fixed decoding
method, and a rule that extracts the answer from the completion and compares it with the gold answer. In this
problem you run such a benchmark yourself, and you see that the prompt is part of the measurement.

We use GSM8K~\cite{gsm8k}, a set of grade-school math word problems whose answers are numbers, and evaluate on
its first 100 test problems (\file{data/gsm8k\_first\_100.jsonl}). The two models are released instruct models
that you do not train: \textbf{model A} is Qwen2.5-0.5B-Instruct and \textbf{model B} is
Qwen2.5-1.5B-Instruct (both are downloaded from Hugging Face the first time they are used). This problem is
evaluation only: there is no training, and the whole job takes about 16 minutes on one GPU, mostly spent loading the models.

\file{posttrain/gsm8k.py} provides the benchmark: \texttt{standard\_prompt\_template()} (used in
part~1), \texttt{extract\_answer()} and \texttt{is\_correct()} (numeric comparison with the gold answer).
Like most GSM8K evaluation harnesses, the extractor is lenient: it reads the first number inside the last
\texttt{<answer>}\ldots\texttt{</answer>} block; if there is none, the first number inside the last
\verb|\boxed{...}|; and if there is neither, the last number anywhere in the completion (commas are removed
first). For this problem the \emph{format rate} is the fraction of completions that give an explicit final
answer, in \texttt{<answer>} tags or \verb|\boxed{}|. \file{scripts/eval\_gsm8k.py} renders each prompt with
the model's chat template, decodes greedily with at most 512 new tokens, stops after the first
\texttt{</answer>}, and writes \file{predictions.jsonl} and \file{metrics.json} (accuracy and format rate).
\begin{enumerate}
  \item (Experiment, 2 pts) Evaluate both models with the standard prompt:
\begin{lstlisting}
PROMPTS=standard sbatch slurm/gsm8k.sbatch   # -> results/gsm8k/{A_standard,B_standard}/\end{lstlisting}
        Each run earns 1 point if its \file{predictions.jsonl} contains all 100 problems and every prompt in it
        is the standard prompt. There is no accuracy threshold.
  \item (Programming and experiment, 1 pt) Write \texttt{improved\_prompt\_template()} in
        \file{posttrain/gsm8k.py}: a prompt that makes \textbf{model A} more accurate than the standard prompt
        does. You may change the instructions, add a system message, or add few-shot example turns (take
        examples from the GSM8K \emph{training} set or write your own, never from the 100 evaluation problems).
        The model, the decoding settings and the answer extractor stay fixed. Ask for the final answer as
        \texttt{<answer>NUMBER</answer>} so that the extractor reads the right number. Then run
\begin{lstlisting}
PROMPTS=improved sbatch slurm/gsm8k.sbatch   # -> results/gsm8k/A_improved/\end{lstlisting}
        You earn the point if the autograder's re-scored accuracy of \file{A\_improved} is strictly higher than
        that of your \file{A\_standard}.
  \item (Write-up, 2 pts) Include \file{results/gsm8k/bar.png} (made by \file{plot\_results.py}) and report
        the accuracy and format rate of the three runs. Show your prompt and explain why it helps, using what you
        see in the \file{A\_standard} and \file{A\_improved} completions. Compare model A with your prompt to
        model B with the standard prompt: which model looks better? What does this mean for a paper that
        compares two models using a single fixed prompt?
\end{enumerate}
\begin{solution}
\resultfigure{0.6\linewidth}{results/gsm8k/bar.png}{GSM8K accuracy on 100 test problems.}

The accuracy and format rate were [accuracy, format rate] for
\texttt{A\_standard}, [accuracy, format rate] for \texttt{B\_standard},
and [accuracy, format rate] for \texttt{A\_improved}.

\begin{lstlisting}
System
Solve the word problem with short calculations, then put only the final
number inside <answer>NUMBER</answer>.

User
A kitchen filled 6 trays with 8 rolls each, then sold 19 rolls.
How many rolls were left?

Assistant
6 * 8 = 48 rolls. 48 - 19 = 29 rolls left.
<answer>29</answer>

User
A shop sold 9 notebooks for 4 dollars each and 3 folders for
2 dollars each. How many dollars did the shop collect?

Assistant
9 * 4 = 36 and 3 * 2 = 6. 36 + 6 = 42 dollars.
<answer>42</answer>

User
[The evaluation question.]
\end{lstlisting}

The prompt gives the model space to write intermediate calculations and
demonstrates where the final number belongs. On [same test ID], the
standard prompt produced [observed completion and extracted answer],
while the revised prompt produced [observed completion and extracted
answer]. This example suggests [explanation supported by the outputs].

Under the prompts tested, [model and prompt] scored higher by
[difference] percentage points. A comparison with one fixed prompt
can reflect how well each model responds to that prompt, so the ranking
may change under a different reasonable prompt.
\end{solution}

% =====================================================================================
\section*{Disclosure (required)}
Describe how you used AI tools for this assignment (see ``Use of AI Tools'' at the beginning). For each tool,
give its name (for example ChatGPT, Claude, Gemini or GitHub Copilot), what you used it for (for example
explaining a concept, explaining an error message, debugging a specific function, or editing the wording of the
write-up), and which problems or files it touched. If you did not use any AI tools, write ``I did not use AI
tools for this assignment.''
\begin{solution}
I used OpenAI Codex to summarize the assignment and propose and apply a few-shot
GSM8K prompt for Problem 7.
\end{solution}

% =====================================================================================
\section*{Grading and Code Expectations}
\begin{itemize}
  \item \textbf{Unit tests (50 pts)} run against \textbf{hidden} test cases that check the same properties as the
        public tests but use different numbers. Points per test file are proportional to the fraction of its
        tests that pass (common 12, sft 6, preference 10, grpo 14, test\_time 8). Passing the public tests is
        necessary but not sufficient, so do not hard-code values.
  \item \textbf{Results (26 pts).} The autograder re-scores your \file{predictions.jsonl} with its own verifier
        (\file{metrics.json} is not trusted). You get full credit at or above the threshold and partial credit
        proportional to accuracy below it. Thresholds are 90\% of the staff reference accuracy.
  \item \textbf{GSM8K (3 pts).} The autograder re-scores your three GSM8K \file{predictions.jsonl} files with its own
        copy of \texttt{extract\_answer} and \texttt{is\_correct}; see Problem 7 for what earns each point.
  \item \textbf{What Gradescope shows.} Right after you submit you see whether all files were found, your
        re-scored SFT/DPO/GRPO accuracies and points, the test-time artifact check, the GSM8K points, and the public unit tests (worth
        0 points; the same tests as in your starter code). The hidden unit tests and your total score appear when
        grades are published. You may resubmit until the deadline; the last submission counts.
  \item \textbf{Academic integrity.} You may use any resources to understand the material, but you may not copy
        code from online resources (including public solutions to similar assignments); write your own code.
        AI tools may be used only as described in ``Use of AI Tools'', and every use must be reported in the
        Disclosure section.
  \item \textbf{Do not edit} \file{countdown.py}, \file{evaluate.py}, \file{generation.py},
        \file{eval\_gsm8k.py}, the provided functions in \file{gsm8k.py}, or the decoding settings. Predictions must come from \file{scripts/evaluate.py} run on the model that the corresponding
        training script produced; for example, selecting among several samples with the verifier is not allowed.
        Keep your \file{results/*/final} checkpoints on Great Lakes until grades are released: we spot-check
        submissions by re-running the evaluation on them. Edited or substituted predictions are an honor-code
        violation.
  \item \textbf{Write-up (17 pts)} is graded by hand. A question gets full credit when the requested plot or
        number is present and the explanation is correct, and half credit when only one of the two is. Plots need
        axis labels, and claims contradicted by your own numbers get no credit.
\end{itemize}

% =====================================================================================
\section*{Computation Guide}
NLP experiments are commonly performed in a Linux-based shell environment. We recommend using
\href{https://arc.umich.edu/greatlakes/}{Great Lakes}, a Linux-based computing cluster at UMich, for this
assignment; the provided Slurm scripts target its \texttt{spgpu} partition (one NVIDIA A40, 48\,GB).
\textbf{If you have not already, please sign up for a Great Lakes account as soon as possible.}

\noindent\textcolor{red}{\textbf{Please DO NOT use this account for anything other than coursework, unless given
permission from the instructors. Failure to comply with these requirements will lead to a zero score for this
assignment.}}

\begin{itemize}
  \item \textbf{VPN.} To access Great Lakes when not on campus, you will need to use the UMich VPN.
  \item \textbf{Great Lakes dashboard.} The online dashboard
        (\href{https://greatlakes.arc-ts.umich.edu}{greatlakes.arc-ts.umich.edu}) offers a file manager
        (\texttt{Files > Home Directory}), a terminal (\texttt{Clusters > \_greatlakes Shell Access}) and a job
        composer. You can also SSH into the cluster from a local terminal.
  \item \textbf{Slurm.} You \textbf{MUST} use Slurm jobs (\texttt{sbatch}) for training and evaluation; the cluster
        kills long or memory-hungry processes on login nodes. You do not need to edit the provided
        \file{slurm/*.sbatch} files: they already set the course account (\texttt{cse595f26\_class}), the partition
        and the resources, and \file{slurm/env.sh} sets up the environment. Always run \texttt{sbatch} from the top
        of your starter-code folder. The course account allows each student \textbf{one GPU job at a time} and at
        most \textbf{4 hours per job}, and the whole class shares 20 GPUs, so jobs may wait in the queue: submit
        the next stage as soon as the previous one finishes.
  \item \textbf{Python environment.} \file{slurm/setup\_env.sh} creates a virtual environment in your scratch
        space with the pinned package versions in \file{requirements.txt}. Model downloads go to scratch, not
        your home directory.
\end{itemize}

\paragraph{Compute budget.} Measured wall times of the staff reference run on one A40:
\begin{center}
\begin{tabular}{@{}lp{6.6cm}r@{}}
\toprule
job & what it runs & wall time \\
\midrule
\file{sft.sbatch} & base-model evaluation, SFT (8000 examples $\times$ 2 epochs), evaluation & $\approx$ 20 min \\
\file{dpo.sbatch} & DPO on the provided pairs, evaluation & $\approx$ 10 min \\
\file{grpo.sbatch} & GRPO (100 steps), evaluation & $\approx$ 70 min \\
\file{grpo\_ablation.sbatch} & GRPO ablation (40 steps) & $\approx$ 30 min \\
\file{test\_time.sbatch} & 16 samples $\times$ 200 problems, GRPO and SFT models & $\approx$ 10 min \\
\file{gsm8k.sbatch} & 3 GSM8K evaluations (100 problems each), no training & $\approx$ 16 min \\
\midrule
total & & $\approx$ 2.6 h \\
\bottomrule
\end{tabular}
\end{center}
Please smoke-test with small sizes first: a crashed long job wastes your time and the shared queue. The default
settings need a 48\,GB GPU such as the A40. If you hit out-of-memory errors, lower \texttt{-{}-micro-batch-size} or
add \texttt{-{}-grad-checkpointing}; neither changes the math. With \texttt{-{}-micro-batch-size 4 -{}-grad-checkpointing}
(SFT and GRPO) and \texttt{-{}-micro-batch-size 2 -{}-grad-checkpointing} (DPO) every stage fits in 16\,GB. The V100
GPUs in the \texttt{gpu} partition are too old for the PyTorch version we use, so do not submit there.

\section*{Helpful Hints and Resources}
\begin{itemize}
  \item In \texttt{grpo\_microbatch\_train\_step}, divide by \texttt{gradient\_accumulation\_steps}
        \textbf{before} calling \texttt{backward()}.
  \item A flat GRPO reward with \texttt{frac\_groups\_all\_same} $\approx 1$ means there is no learning signal.
        Check your SFT model's accuracy first.
  \item To see what your model actually generates, read the \texttt{completion} field of each line of
        your \file{predictions.jsonl} files.
  \item The reference papers below give the derivations behind each loss; the docstrings give the exact
        conventions the tests expect.
\end{itemize}

\begin{thebibliography}{99}
\bibitem{qwen25} Qwen Team. ``Qwen2.5 Technical Report.'' arXiv preprint arXiv:2412.15115 (2024).
\bibitem{dpo} Rafailov, R., Sharma, A., Mitchell, E., Ermon, S., Manning, C. D., \& Finn, C. ``Direct Preference
  Optimization: Your Language Model is Secretly a Reward Model.'' Advances in Neural Information Processing Systems (2023).
\bibitem{ipo} Azar, M. G., Rowland, M., Piot, B., Guo, D., Calandriello, D., Valko, M., \& Munos, R. ``A General
  Theoretical Paradigm to Understand Learning from Human Preferences.'' In International Conference on Artificial
  Intelligence and Statistics (2024).
\bibitem{instructgpt} Ouyang, L., Wu, J., Jiang, X., Almeida, D., Wainwright, C. L., Mishkin, P., et al. ``Training
  Language Models to Follow Instructions with Human Feedback.'' Advances in Neural Information Processing Systems (2022).
\bibitem{deepseekmath} Shao, Z., Wang, P., Zhu, Q., Xu, R., Song, J., Bi, X., et al. ``DeepSeekMath: Pushing the
  Limits of Mathematical Reasoning in Open Language Models.'' arXiv preprint arXiv:2402.03300 (2024).
\bibitem{selfconsistency} Wang, X., Wei, J., Schuurmans, D., Le, Q., Chi, E., Narang, S., Chowdhery, A., \& Zhou, D.
  ``Self-Consistency Improves Chain of Thought Reasoning in Language Models.'' In International Conference on
  Learning Representations (2023).
\bibitem{codex} Chen, M., Tworek, J., Jun, H., Yuan, Q., Pinto, H. P. d. O., Kaplan, J., et al. ``Evaluating Large
  Language Models Trained on Code.'' arXiv preprint arXiv:2107.03374 (2021).
\bibitem{streamofsearch} Gandhi, K., Lee, D., Grand, G., Liu, M., Cheng, W., Sharma, A., \& Goodman, N. D. ``Stream
  of Search (SoS): Learning to Search in Language.'' In Conference on Language Modeling (2024).
\bibitem{rloo} Ahmadian, A., Cremer, C., Gall\'e, M., Fadaee, M., Kreutzer, J., Pietquin, O., \"Ust\"un, A., \&
  Hooker, S. ``Back to Basics: Revisiting REINFORCE-Style Optimization for Learning from Human Feedback in LLMs.''
  In Proceedings of the Association for Computational Linguistics (2024).
\bibitem{ppo} Schulman, J., Wolski, F., Dhariwal, P., Radford, A., \& Klimov, O. ``Proximal Policy Optimization
  Algorithms.'' arXiv preprint arXiv:1707.06347 (2017).
\bibitem{dapo} Yu, Q., Zhang, Z., Zhu, R., Yuan, Y., Zuo, X., Yue, Y., et al. ``DAPO: An Open-Source LLM
  Reinforcement Learning System at Scale.'' arXiv preprint arXiv:2503.14476 (2025).
\bibitem{drgrpo} Liu, Z., Chen, C., Li, W., Qi, P., Pang, T., Du, C., Lee, W. S., \& Lin, M. ``Understanding
  R1-Zero-Like Training: A Critical Perspective.'' arXiv preprint arXiv:2503.20783 (2025).
\bibitem{gsm8k} Cobbe, K., Kosaraju, V., Bavarian, M., Chen, M., Jun, H., Kaiser, L., et al. ``Training Verifiers
to Solve Math Word Problems.'' arXiv preprint arXiv:2110.14168 (2021).
\end{thebibliography}

\end{document}
